\section{Préambule}
Ce document décrit l'installation, l'exploitation et la résolution de pannes associées au composant "Pare-feu externe" de la solution "Infr@home".\\

\subsection{Le projet PFsense}
PFsense est une solution de pare-feu s'appuyant sur le système d'exploitation FreeBSD. Le suivi du projet est assuré par la société Netgate. Le site \url{https://www.pfsense.org/} met à disposition le téléchargement des sources et de nombreuses documentations relatives au projet.

\begin{warning}
    Ce pare-feu est choisi dans le cadre du projet car il est disponible en sources ouvertes. Cependant, les deux pares-feux ne devraient pas être issus d'un même développement ; un pare-feu au moins devrait être un produit qualifié par l'ANSSI (idéalement le plus proche d'Internet).
\end{warning}

\subsection{Configuration requise}
La configuration minimale \footnote{\url{https://docs.netgate.com/pfsense/en/latest/book/hardware/minimum-hardware-requirements.html}} (trafic non-chiffré à 100Mo/s) pour PFsense est définie ci-dessous :
\begin{itemize}
    \item \textbf{Processeur} : 1 CPU monocore
    \item \textbf{Mémoire} : 1 Go de RAM
    \item \textbf{Stockage} : 20 Go d'espace disque
    \item \textbf{Réseau} : Au moins une interface par réseau soutenu
\end{itemize}

% Les performances requises peuvent être affinées en fonction des fonctionnalités à activer. Le détail des calculs est donné sur le site de PFsense\footnote{\url{https://docs.netgate.com/pfsense/en/latest/book/hardware/hardware-sizing-guidance.html#feature-considerations}}.
Les performances requises peuvent être affinées en fonction des fonctionnalités à activer. Le détail des calculs est donné sur le site de PFsense\footnote{\url{https://docs.netgate.com/pfsense/en/latest/book/hardware/hardware-sizing-guidance.html\#feature-considerations}}.\\
~\\
Il est ensuite nécessaire de récupérer la dernière image ISO de PFSense depuis le site de l'éditeur, en architecture x64 : \url{https://www.pfsense.org/download/}.\\
~\\
L'installation peut alors commencer.